\cotefile{81}
\scanpage{27}
% loc: 1–3
cette dernière relation $\Sigma_1^-\cap\Sigma_2^-\ne\varnothing$ est vraie \unsure{=|même} si $E_1=E_2=\varnothing$, car \deleted{\unsure{alors par hyp.|par hyp.}} si on avait de plus $\Sigma_1^-\cap\Sigma_2^-=\varnothing$, on en
% loc: 4–8
conclurait que les deux \deleted{\unsure{part.}} décompositions binaires \deleted{\unsure{son|se}} (qui sont alors des partitions binaires) sont \uline{identiques}, ce que l'on \inserted{\unsure{a|à}} exclu. Cela montre que les $\Sigma_1^+,\Sigma_2^+$ sont uniquement déterminés.

% loc: 9
On a donc le diagramme d'\unsure{inclusions|inclusion}
% loc: upper chain
\diagram{\begin{tikzcd}[ampersand replacement=\&,column sep=small]
\Sigma_1^+ \& E_1 \arrow[l,hook] \arrow[r,hook] \& \Sigma_{12} \& E_2 \arrow[l,hook] \arrow[r,hook] \& \Sigma_2^+
\end{tikzcd}}
% loc: 10
\unsure{qu'on|qui est|qu'on a}, \unsure{exprime|exprimé} par \unsure{nos diagrammes|un diagramme}.
% loc: lower chain
\diagram{\begin{tikzcd}[ampersand replacement=\&,column sep=small,row sep=2pt]
\unsure{\bullet|\circ} \arrow[rr,dash] \& {} \& \unsure{\bullet|\blacktriangledown} \arrow[rr,dash] \& {} \& \bullet\\
\Sigma_1^+ \& E_1 \arrow[l,hook] \arrow[r,hook] \& \Sigma_{12} \& E_2 \arrow[l,hook] \arrow[r,hook] \& \Sigma_2^+
\end{tikzcd}}
% loc: NB at right of lower chain
\margin{\[
\text{NB}\left\{\begin{aligned}
\Sigma_1^-&=\Sigma_{12}\amalg_{E_2}\Sigma_2^+\\
\Sigma_{\deleted{\unsure{1}}2}^-&=\deleted{\unsure{E}}\Sigma_{12}\amalg_{E_1}\Sigma_1^+
\end{aligned}\right.
\]}
% slanted
% loc: slanted note, lines 1–8
\margin{Ceci n'implique pas que les quatre ensembles $\Sigma_1^+,\Sigma_1^-,\Sigma_2^+,\Sigma_2^-$ soient distincts, car on peut avoir (\unsure{non pas|bien que|bien pas} $\Sigma_1^+=\Sigma_2^+$, ni $\Sigma_1^-=\Sigma_2^-$, mais) $\Sigma_1^+=\Sigma_2^-$, qui équivaut à $E_1=\Sigma_{12}$ et entraîne $E_2=\varnothing$
% loc: slanted note, lines 9–13
(Cela signifie, en termes des $(\Sigma_{\unsure{1|t|i}}^+,\Sigma_1^-)$ que $\Sigma_2^{\unsure{\dot{-}|-}}=\Sigma_1^+$, $\Sigma_{\unsure{2|t}}^+=\unsure{X|\varnothing}\setminus\Sigma_1^+=\overset{\unsure{\circ|\leftarrow}}{\Sigma}{}_1^{-,\unsure{\circ|0}}$) \deleted{\unsure{ou}} et la relation \unsure{symetrique|symétrique}. On ne veut \unsure{pas exclure} \unsure{ce|ces} cas, \unsure{i.e.|ici} \deleted{\unsure{mais|nous}} n'\unsure{exige pas|exigeons pas} que
% loc: slanted note, lines 14–19
les quatre ensembles $\Sigma_1^+,\Sigma_1^-,\Sigma_2^+,\Sigma_2^-$ soient distincts \unsure{s'ils|ssi ils} le \unsure{\deleted{sont}|sont}, \unsure{on|c.-à-d.} \unsure{dit|dire} que les \deleted{\unsure{ens.}} \inserted{deux \unsure{décomp. bin.|décomp. bis}} \deleted{\unsure{sont strictmt}} \unsure{non parallèles|sont non parallèles}.}
% loc: 11–12, main text below chains
Je me propose d'étudier \unsure{les|des} ensembles finis \unsure{indicellés|indiciellés|indécollés|indécollt} $\deleted{\unsure{\Pi|\Sigma}}(\pi_\alpha)_{\alpha\in A}\deleted{\unsure{=}(\unsure{\Sigma_\alpha^i}),\unsure{\alpha\in A}}$
% loc: 13–14
$\deleted{\pi_\alpha=\{\Sigma_\alpha^\varepsilon}$ \unsure{de|des} décompositions binaires de $\unsure{X}$ deux à deux non parallèles.

% loc: 15–16
Voici une façon \unsure{très} \unsure{générale|génerale} d'\deleted{\unsure{en}}obtenir des \unsure{floppées|floppés}.
% archive sheet number: 26
\scanpage{28}
% loc: 1–3
Soit \unsure{.} $\Gamma=(S,A,R)$ un arbre ayant $S$ comme \unsure{ens|ens.} de sommets, $A$ \unsure{-} comme ens d'\unsure{arêtes|aretes}, $R$ comme \unsure{rélation|relation|rèlation} d'incidence \unsure{[|E}\unsure{donc|d'où}
% loc: 4–6
$A\hookrightarrow\unsure{\mathcal{P}|\mathscr{P}}_2(S)$, et $S\unsure{\hookrightarrow|\leftrightarrow|\mapsto}\mathcal{P}(A)$, et cette application est injective sauf si $\Gamma$ est \unsure{isom|isom.} au graphe
\diagram{\begin{tikzcd}[ampersand replacement=\&,column sep=1.4em]
\bullet \arrow[r,dash] \& \bullet
\end{tikzcd}}
]. \deleted{\unsure{Soit}}Soient $(U_s)_{s\in S}$ et $(E_a)_{a\in A}$ deux familles d'ensembles, et posons
% loc: display below families
\[
X=\coprod_{s\in S}U_s\ \amalg\ \coprod_{a\in A}E_a
\]

% loc: 9–12
Pour tout $a_0\in A$, considérons les deux \unsure{comp|comp.} \unsure{conn.|conn} $\deleted{\unsure{E}}\Gamma',\Gamma''$ de l'arbre $\Gamma$ ``coupé en $a_0$'', et \deleted{\unsure{con|c}}considérons la décomposition binaire \unsure{.} $X=X'\cup X''$ de $X$,
% loc: 13 and two-line display
où
\[
\begin{aligned}
X'&=\coprod_{s\in S'}U_s\amalg\coprod_{a\in A'}E_a\\
X''&=\coprod_{s\in S''}U_s\amalg\deleted{\unsure{E}}\coprod_{\unsure{a|c}\in A''}E_a
\end{aligned}
\]
% loc: intersection
avec
\[
X'\cap X''=E_{a_0}
\]
% loc: 16–20
cette décomposition est ``propre''\ \unsure{sss|ssi} $\deleted{\unsure{E|\in}}X'\ne E_{a_{\unsure{0|b}}},\ X''\ne E_{a_0}$. Ceci est le cas, quel que soit le choix de $a_0\in A$, sss on \unsure{a|à} la propriété \unsure{suivante|suivant} :

% loc: bottom condition
$(\unsure{\ast|A})$ Si $s\in S$ est un sommet-bout, alors $U_s\ne\varnothing$
% archive sheet number: 27
\scanpage{29}
% slanted
% loc: upper left slant
\margin{$(\unsure{\ast\ast|\ast})$ \unsure{exprime|expriment|exprimons|exigeons} que si $a\ne b$, alors $\pi_a\ne\pi_b$}
% loc: 1–4
De plus, si $a_1,a_2\in A$, $a_1\ne a_2$, on voit de suite que les décompositions $\pi_{a_1},\pi_{a_2}$ sont bien parallèles, à cela près qu'il faut exiger \deleted{\unsure{que si $s$ est un sommet d'or}}%
% loc: line 4, below deletion and continued in left margin
\inserted{$(\unsure{\ast\ast})$ si $s\in S$ est d'ordre $2$, alors un des trois \unsure{ens|ens.} $U_s,E_a,E_b$ est $\ne\varnothing$, où $a,b$ sont les deux arêtes \unsure{inc.|inci.} à $s$}

% loc: 5–8
Une \unsure{\deleted{façon}|façon} de visualiser la construction est via un arbre topologique $\unsure{\mathcal{X}|\mathscr{X}|X}\deleted{\unsure{A|\Gamma}}=(\overset{\text{\inserted{\unsure{esp. top.}}}}{\underline{\Gamma}},S)$ (tel que $\Gamma$ en soit une description combinatoire). On fixe un pt sur
% loc: 9–12
chaque \unsure{arête|arète}, de façon à \unsure{donner|trouver} un graphe topologique (subdivision barycentrique) $\mathcal{X}^{\ast}=(\underline{\Gamma}^{\unsure{\ast}},\underbrace{S\amalg A'}_{S^{\ast}})$, où $A'\unsure{\subset|\in}\underline{\Gamma}$ est en \unsure{corr|corr.} \deleted{\unsure{b|bi}}biunivoque donnée avec \deleted{\unsure{tout}} $A$.

% loc: 13–17, bracketed NB
\deleted{\unsure{pt de}} [\uline{NB} $A'$ est formé de \deleted{pts} sommets d'ordre $2$, et \inserted{pour} \unsure{tte|toute} arête de $\mathcal{X}^{\ast}$, une exactement des deux extrémités appartient à $A'$, l'autre à $S$ -- voici une description complète des propriétés particulières \unsure{aux|des} $\underline{\Gamma}^{\ast}$\unsure{..|...}]

% loc: 18–21
Alors \deleted{les}\unsure{le} point\deleted{s} \inserted{$a'$} de $A'$ coupe\deleted{\unsure{nt}} $\underline{\Gamma}^{\deleted{\unsure{\ast}}}$ en deux composantes connexes, dont les adhérences définissent une décomposition binaire de $\underline{\Gamma}$, d'\unsure{``|''}équateur $\{a'\}$. Deux
% loc: 22–23
telles décompositions \deleted{\unsure{bi}} de $\underline{\Gamma}$, pour $a'\ne\unsure{a''|a'}$, définissent des décompositions binaires \unsure{strict$^t$|stric$^t$} non parallèles.
% archive sheet number: 28
\scanpage{30}
% loc: 1–4
\unsure{Item|Idem|Itou}, si on regarde les traces de ces décompositions binaires sur $S^\ast=S\amalg A'$. Donc pour \unsure{tte|toute} application \uline{surjective} $X\unsure{\mathrel{\text{---}}|\longrightarrow}S^\ast$, si on regarde \unsure{les des|les les|là les|les} images
% loc: 5–8
inverses des décompositions binaires précédentes de $S^\ast$. \deleted{\unsure{D}}Or la donnée d'une telle application\unsure{,} revient à celle d'une famille $(X_\sigma)_{\sigma\in S^{\unsure{\ast|+}}}$ \inserted{d'ens $\ne\varnothing$}, i.e. \deleted{\unsure{à}} celle
% loc: 9–10
\unsure{des|de} deux familles $(U_s)_{s\in S}$ et $(E_a)_{a\in A}$ d'ens non vides.

% loc: 11–15
Le cas où certains des ensembles $U_s,E_a$ pourraient être vides est à examiner avec plus de soin, il donne néanmoins ce qu'on veut \unsure{sss|ssi} on a les conditions $(\ast)$ $(\ast\ast)$ ci-dessus, et la condition de
% loc: 16–17
non-parallélisme stricte sss on a \uline{ceci} (plus fort que $(\unsure{\ast\ast|\ast})$)
% loc: 18–20, condition at foot
$(\ast\ast\ast)$ Si $s\in S$ est un sommet \uline{d'ordre 2} \deleted{\unsure{tel}} \inserted{et} \deleted{que $U_s=\varnothing$, et} \inserted{$a\in A$ une arête incidente, on a}
\[
U_s\ne\varnothing\ \unsure{\text{ou}|\text{--}}\ E_a\ne\varnothing
\]
% archive sheet number: 29
\scanpage{32}
% archive sheet number: 31
\begingroup
\makeatletter\tagsleft@true\makeatother
% loc: Upper left and title
% slanted
\margin{\uline{Version II}}
\noindent Décompositions $n$-aires d'un espace topologique\unsure{.}

% loc: 1.1, opening and pair
1.1. Soit $X$ un espace topologique. On appelle \uline{décomposition binaire} de $X$ une paire (\unsure{notée|noté} \unsure{indiciellement})
\[
\pi=(\Sigma^\varepsilon)_{\varepsilon\in\omega}\qquad(\omega\ \text{ens à deux éléments})
\]
% loc: 1.1, before (1)
\unsure{dé|de|d'e} \deleted{ensembles} parties \uline{fermées} (distinctes) de $X$, de réunion $X$
\[
X=\bigcup_{\varepsilon\in\omega}\Sigma^\varepsilon.\tag*{(1)}
\]

% loc: (2) and (2')
On suppose que
\[
\Sigma^\varepsilon\ne\varnothing\qquad(\forall\varepsilon\in\omega)\tag*{(2)}
\]
ce qui signifie aussi que
\[
\pi\ne\{\varnothing,X\}\tag*{$(2')$}
\]
% loc: condition before (3) and (3)
et implique que $X\ne\varnothing$, et même $\mathrm{card}\,X\geq 2$, puisque l'on a supposé
\[
\deleted{\unsure{X|\Sigma}}\Sigma^\varepsilon\deleted{\unsure{=|\ne}}\ne\Sigma^{\varepsilon'}\qquad\text{si }\varepsilon,\varepsilon'\in\omega,\quad\varepsilon\ne\unsure{\varepsilon'|r'}.\tag*{(\unsure{3|s|$\wedge$})}
\]

% loc: inclusion and following equality
On n'exclut pas le cas où on aurait une inclusion
\[
\Sigma^{\varepsilon'}\subset\Sigma^\varepsilon
\]
ce qui équivaut, par (1), à la relation
\[
\Sigma^\varepsilon=X.
\]

% loc: hémisphères and (4)
Les $\Sigma^\varepsilon$ $(\varepsilon\in\omega)$ s'appellent les \uline{\unsure{hémisphères|hémisphéres}} de la décomposition binaire, \deleted{\unsure{la|le}} et l'ensemble
\[
E=\bigcap_{\deleted{\unsure{\omega|\varepsilon}}\varepsilon\in\omega}\Sigma^\omega\tag*{(4)}
\]
s'appelle l'\uline{équateur} de la décomposition binaire.

% loc: last paragraph and interline additions
On n'exclut pas le cas où on aurait $E=\varnothing$, ce qui implique que les $\Sigma^\varepsilon$ \inserted{$(\varepsilon\in\omega)$} sont \inserted{complémentaires l'un de l'autre, \unsure{donc}} à la fois fermés et ouverts -- donc que $X$ est disconnexe.
\endgroup
\scanpage{33}
% archive sheet number: 32
\begingroup
\makeatletter\tagsleft@true\makeatother
% loc: opening and (5), upper line
Notons les implications
\[
\left\{\begin{aligned}
&X\ \text{connexe}\quad\Longrightarrow\quad E\overset{\mathrm{def}}{=}\bigcap_{\varepsilon\in\omega}\deleted{\unsure{X|\Sigma}}\Sigma^\varepsilon\ \deleted{\unsure{-|=}}\deleted{\unsure{=|\varnothing}}\ne\varnothing\\
% loc: (5), lower line and left of brace
&\text{\unsure{\deleted{Si}|Si|si} les }\Sigma^\varepsilon\ \text{connexes, alors}\quad X\ \text{connexe}\ \Longleftrightarrow\ E\ne\varnothing.
\end{aligned}\right.\tag*{\shortstack[l]{(5)\\\unsure{,|'}}}
\]

% loc: (6), definition
Nous poserons
\[
\Sigma^{\unsure{\circ|o|0}\varepsilon}=\Sigma^\varepsilon\setminus E=\Sigma^\varepsilon\setminus\Sigma^\varepsilon\cap\Sigma^{\varepsilon'}=X\setminus\Sigma^{\varepsilon'}\qquad(\unsure{\omega|w}\unsure{\ne|=|\in}\{\varepsilon,\varepsilon'\})\tag*{(6)}
\]
% loc: (7), modified exponents and parenthesis
de sorte que l'on a la décomposition
\[
X=\Sigma^{\circ\deleted{\unsure{\varepsilon|2}}\varepsilon_{\deleted{\unsure{1|+}}\unsure{.|,}}}\cup E\cup\Sigma^{\circ\varepsilon'_{\deleted{\unsure{2|-}}}}\qquad(\omega=\{\unsure{\varepsilon'|\varepsilon}_{\deleted{\unsure{1|+}}},\varepsilon'_{\deleted{\unsure{2|-}}}\})\tag*{(7)}
\]
de $X$ en \deleted{trois} ens. mutuellement disjoints.

% loc: (8), bars, point and punctuation
Notons
\[
\overline{\Sigma^{\circ\varepsilon}}\cap E=\overline{\Sigma^{\circ\varepsilon}}\cap\Sigma^{\varepsilon'}=\overline{\Sigma^{\circ\varepsilon}}\setminus\Sigma^{\circ\varepsilon}=\dot{\Sigma}^{\circ,\varepsilon}\quad\unsure{;|\varepsilon,|\varepsilon|,}\tag*{(8)}
\]

% loc: two struck lines before 1.2
\deleted{On a le \unsure{cas}}\par
\deleted{\unsure{on|On} s'intéresse à la \unsure{condi|cond}}

% loc: 1.2 and (9)
1.2. Deux décompositions binaires
\[
\pi_1=\{\Sigma_{\deleted{\unsure{\varepsilon|2}}1}^\varepsilon\}_{\varepsilon\in\omega_1},\qquad\pi_2=\{\Sigma_2^\varepsilon\}_{\varepsilon\in\omega_2}\tag*{(9)}
\]
% loc: line before (10) and insertion
sont dites non \uline{\unsure{parallèles|paralèles}}, \inserted{si elles sont distinctes, et} s'il existe un \unsure{.|,} $\Sigma_1^+\in\pi_1$, et un $\Sigma_2^+\in\pi_2$, tels que l'on ait
\[
\Sigma_1^+\cap\Sigma_2^+=\varnothing\tag*{(10)}
\]
% loc: after (10), (11) and (11 bis)
condition qui équivaut à
\[
\Sigma_1^+\subset\Sigma_2^-\tag*{(11)}
\]
ou encore à
\[
\Sigma_2^+\subset\Sigma_1^-\tag*{(\unsure{11bis|11 bis|11'bis})}
\]
% loc: last line and (12)
où $\Sigma_1^-,\Sigma_{\unsure{2|1}}^-$ sont \unsure{définis|définies} par
\[
\pi_1=\{\Sigma_1^+,\Sigma_1^-\},\qquad\pi_2=\{\Sigma_2^+,\Sigma_{\unsure{2|1}}^-\}\tag*{(12)}
\]
\endgroup
\scanpage{34}
% archive sheet number: 33
\begingroup
\makeatletter\tagsleft@true\makeatother
% loc: opening, (13)
Posons
\[
\Sigma_{12}=\Sigma_1^-\cap\Sigma_2^-,\qquad E_1=\Sigma_1^+\cap\Sigma_1^-,\quad E_2=\Sigma_1^+\cap\Sigma_2^-\tag*{(13)}
\]
alors on a le diagramme d'inclusion de parties fermées de $X$
% loc: diagram (14), lower slant
\usetikzlibrary{decorations.pathmorphing}
\diagram{\[
\begin{tikzcd}[ampersand replacement=\&,row sep=1.1em,column sep=.4em]
\& \& X \& \& \\
\& \Sigma_2^- \arrow[ur] \& \& \Sigma_1^- \arrow[ul] \& \\
\Sigma_1^+ \arrow[ur] \& \& \Sigma_{12} \arrow[ul]\arrow[ur] \& \& \Sigma_2^+ \arrow[ul]\\
\& E_1 \arrow[ul]\arrow[ur] \& \& E_2 \arrow[ul]\arrow[ur] \& \\
\& \& \deleted{\unsure{\Omega|2|\varnothing}}\arrow[ul]\arrow[ul,no head,decorate,decoration={zigzag,segment length=3.4pt,amplitude=1.6pt}] \& \&
\end{tikzcd}\tag*{(14)}
\]}
% loc: right marginal NB alongside (14)
\margin{NB $\Sigma_1^+\cap\Sigma_2^+=\varnothing$ donc $E_1\cap\Sigma_2^+=\varnothing$, $E_2\cap\Sigma_1^-=\varnothing$ i.e. $E_1\subset\Sigma_2^{\unsure{\circ|o|0}-}$, $E_2\subset\Sigma_1^{\circ-}$ et $E_1\cap E_2\unsure{.|-}=\varnothing$}
% loc: line under (14), cancelled starter and next display
où \unsure{les|le} trois carrés sont cartésiens et cocartésiens
\[
\deleted{\unsure{E_1|E_+}}
\]
i.e.
\[
\left|\begin{aligned}
&E_1=\Sigma_1^+\cap\Sigma_{12},\quad\Sigma_2^-=\Sigma_1^+\cup\Sigma_{12}\\
&\deleted{\unsure{E_1|E}}
\end{aligned}\right.
\]
et les deux relations analogues correspondant aux deux autres carrés
% loc: (15), writes over first row and lower X
\[
\left\{\begin{aligned}
&E_2=\Sigma_{\deleted{\unsure{2|1}}12}^{\deleted{\unsure{+|-}}}\cap\Sigma_{\deleted{\unsure{12|1}}2}^+,\quad \Sigma_1^-=\Sigma_{\unsure{22|12}}^{\deleted{\unsure{+|-}}}\cup\Sigma_2^+\\
&\Sigma_{12}\overset{(13)}{\unsure{\equiv|=}}\Sigma_1^-\cap\Sigma_{\unsure{2|1}}^-,\quad \deleted{\unsure{\Sigma|{-}}}X=\Sigma_2^-\cup\Sigma_1^-.
\end{aligned}\right.\tag*{(15)}
\]

% loc: (16), three lines
Notons que
\[
\left\{\begin{aligned}
&\text{a)}\quad\Sigma_1^+\cap\Sigma_2^-=\Sigma_1^+\ne\varnothing\\
&\text{b)}\quad\Sigma_2^+\cap\Sigma_1^-=\Sigma_2^+\ne\varnothing\\
&\text{c)}\quad\Sigma_1^-\cap\Sigma_2^-=\Sigma_{12}\ne\varnothing,
\end{aligned}\right.\tag*{(16)}
\]

% loc: after (16), interline addition
\inserted{\unsure{Prouvons|Prenons}} \unsure{La|la} dernière relation (les deux premières étant évidentes).
% loc: last paragraph; abandoned inline and display formulas
Si on a\unsure{\deleted{vait}|vait|\deleted{ura}} $\Sigma_{12}=\varnothing$, comme $E_1,E_2\subset\Sigma_{12}$, on aurait $E_1=\varnothing$, $E_2=\varnothing$, de sorte que $\pi_1$ et $\pi_2$ sont deux \uline{partitions} de $X$ en des parties non vides, i.e. $\deleted{\Sigma_{\unsure{1|2}}^- = X\setminus{\unsure{\Sigma|{-}}}^-}$ on \unsure{aurait|avait}
\[
\deleted{\unsure{\Sigma_1^+|\Sigma_1^-}}\quad\Sigma_1^-=X\setminus\Sigma_1^+,\quad\Sigma_2^-=X\unsure{\setminus|-}\Sigma_{\unsure{2|1}}^+
\]
\endgroup
